\section{The original joint Hecke ideal and its characteristic-prime support}
\label{sec:joint-support-producer}

\status{JSP1--JSP12 reconstruct the original ideal-kernel construction,
its full \texttt{RingProps} and \texttt{Dense} consequences, the original
Taylor--Wiles base on the same module and ideal, and the single-point
\texttt{NiceI} implication from the explicitly retained patching inputs
to \texttt{P1core}. They do not construct the retained ordinary
arithmetic system, its four further Hida endpoints, or its all-level
automorphic modules. The joint producer \texttt{COUNT\_II\_SUPPORT}
therefore remains open. Every assertion below is an implication with
its mathematical antecedents printed; a source theorem name is not
used as its proof. The coefficient prime may be three.}

\subsection{The actual objects and the retained arithmetic data}

Fix one of the two original Links of QX4, at a permitted coefficient
object $c'$, and write $O=c'.O$, $k=O/\mathfrak m_O$ and $n=4$.
Retain its original CM field $L$, setting $s$, level $V$, residual
datum $r$, and the already fixed identification
$\iota:\overline{\mathbb Q}_\ell\simeq\mathbb C$.
Write $G_0=V.\mathrm{frame}.G_0$, $G_v$ and $I_v$ for its exact
local groups, and $I_u,I_0,D_0$ for its tame and auxiliary subgroups.
Let $B$ be any original universal ordinary determinant ring, with
pair $(D_B,\psi_B)$, and let $\Lambda$ be any original weight object.
Set $A=\Lambda.\Lambda$ and
\[
 \lambda=\Lambda.\mathrm{toPair}(B.u):A\longrightarrow B.
 \tag{JSP.1}
\]
KX1--KX9 already supply these universal objects and their maps, with
the nonempty-pair antecedent supplied here by QX5.

The ordinary construction uses the original data
\[
 Q=\mathrm{QData.empty}(V.\mathrm{hV}),\qquad
 R:\mathrm{UnitaryHeckeData}(Q,\iota,\Lambda).
\]
Such an $R$ includes
the integral residual type $\tau_0$, the character $\theta$ with its
\texttt{IsTheta} identity, and the residual Hecke eigenvalues $a$.
Its existence is retained. Put
\[
 H=\mathrm{hset}(Q,R.\tau_0),\qquad
 M=H.\mathrm{MK}(R.\theta,R.a,K),\qquad
 \mathbb T=H.\mathrm{TK}(R.\theta,R.a,K),                \tag{JSP.2}
\]
where $K$ is the original empty auxiliary level. The algebra
$\mathbb T$ is the subalgebra of $\operatorname{End}_A(M)$ generated
by the actual Hecke operators. The zero-module case is treated
separately and is never excluded by fiat.

For $M\ne0$ we retain the following ordinary arithmetic input,
on these precise objects.
\begin{enumerate}
\item The Hecke operators commute, $M$ is finite free over $A$,
and $\mathbb T$ is an object of $\mathcal C_O$. Its generators have
the original residual values. These are the relevant control and
integral module assertions, not just a generic existence of forms.
We retain the original $H.\mathrm{Ax}$, in particular its
prime-to-$\ell$ Hecke places and validity of its specified Artin maps.
\item There are an adically continuous determinant $D_{\mathbb T}$
and ordered local characters $\psi_{\mathbb T,v,i}$ with the
literal \texttt{IsUniv} Frobenius, unit and uniformizer identities.
Their reductions are $r.\mathrm{Dbar}$ and $r.\mathrm{psib}$;
the characters have the prescribed torsion restrictions, and
their restrictions to $I_v$ are the images of $\Lambda.\kappa_{v,i}$.
\item At each admitted weight, the original good arithmetic
homomorphisms $\mathbb T\to\overline{\mathbb Q}_\ell$ separate
$\mathbb T$. Their specialized determinants have the original
continuous representations with ordered flags, polarization and
the residual tame type. For the auxiliary assertion in JSP6,
the representations at one such separating weight are trivial
on $I_0$ when $r.\mathrm{V0ok}$ holds.
\end{enumerate}
These hypotheses are the precise parts of the ordinary control,
\texttt{coreDet}, \texttt{Out}, residual-character and auxiliary
point producers that are used below. They are not asserted to have
been supplied by this chapter. In particular, separation at one
weight suffices for the determinant relations, whereas the full
\texttt{Dense} field needs separation at every weight used in JSP7.
For the literal weight dictionary we also retain the following exact
reciprocity matching: every map satisfying the source
\texttt{IsLocalArtin} agrees with its specified $\operatorname{artOf}$,
and on inertia this map satisfies
$\epsilon(\sigma)=\prod_\tau\tau(\operatorname{Art}(\sigma))$.
These source-facing uniqueness and normalization assertions are
distinct from KX3's already proved abstract \texttt{LocCFT} interface.
The source obtains them using global reciprocity, injectivity on
local ideles, and the cyclotomic norm formula. Their exact classical
matching is retained here, not silently attributed to KX3.
The exact source declarations are recorded in the accompanying
provenance file.

% BEGIN JSP1
\paragraph{JSP1: weight coordinates on the original inertia groups.}
\label{claim:JSP1}
Let $d=4[L^+:\mathbb Q]$. There are the source's
$\mathrm{WtVars}(\mathrm{locOf}(r),O,4,d)$ and an isomorphism
\[
 e:O[[X_1,\ldots,X_d]]\xrightarrow{\sim}A
\]
such that, for every coordinate $j$, its chosen
$\gamma_j\in I_{v_j}$, index $i_j$ and unit $c_j\in O^\times$
satisfy
\[
 \overline c_j\,\overline\psi_{v_j,i_j}(\gamma_j)=1,
 \qquad
 e(X_j)=c_j\Lambda.\kappa_{v_j,i_j}(\gamma_j)-1.       \tag{JSP.3}
\]
\begin{proof}
Use KX3's local abelian coordinates on the actual $G_v$ and its
actual inertia, and KX4's fixed character lift
$\psi^0_{v,i}:G_v\to O^\times$. For each basis vector of the
torsion-free inertia factor $\mathbb Z_\ell^{d_v}$, choose a
preimage $\gamma_{v,j}$ in $I_v$. The variable of the universal
principal-unit character at this basis vector is $X_{v,i,j}$.
The universal character therefore has value
$\psi^0_{v,i}(\gamma_{v,j})(1+X_{v,i,j})$ there.
KX5's universal isomorphism to the given $\Lambda$ gives
(JSP.3) with $c_j=\psi^0_{v,i}(\gamma_j)^{-1}$.
Reduction gives its first identity. The number of coordinates is
$\sum_v4d_v=4[L^+:\mathbb Q]$ by KX3.
Enumerate these finitely many triples to obtain the source's
$\mathrm{Fin}(d)$ coordinates; its $\mathrm{Empty}\sqcup\mathrm{Fin}(d)$
index set is the same set with an empty summand. No different local
group, torsion character, or coefficient ring has been substituted.
\end{proof}
% END JSP1

% BEGIN JSP2
\paragraph{JSP2: detecting full determinant identities by arithmetic fibers.}
\label{claim:JSP2}
Let $C$ be a commutative ring and $x_\alpha:C\to E_\alpha$ a
family of homomorphisms which jointly detects zero. Equalities of
determinant laws over $C$, and the source's full determinant-kernel
assertions, can be checked after every $x_\alpha$.
For a determinant law $D$ the latter assertion is literally
\[
 D\text{ kills }z
 \quad\Longleftrightarrow\quad
 D_S(1+h z_S)=1
 \quad\text{for every }C\text{-algebra }S\text{ and }h\in S[G_0].
 \tag{JSP.4}
\]
In a family of polarized flagged representations, the ordered
ordinary product and its star have zero image, and so are killed
in the sense of (JSP.4).
\begin{proof}
For a finite collection of group elements, introduce one independent
variable for each coefficient. Evaluate both determinant laws on
the resulting universal group-algebra element. Their values are
polynomials over $C$. A coefficient which maps to zero under all
$x_\alpha$ is zero; hence the polynomials agree. Naturality then
permits substitution of arbitrary coefficients in any $S$, proving
the asserted law identity. For (JSP.4), use the universal expression
$1+(\sum_j X_j[g_j])z$ and the same coefficient argument. This
does not assume that an injection into a product of fields remains
injective after arbitrary tensor products.

For the local ordinary relation, a common flag
$0=F_0\subset F_1\subset\cdots\subset F_4$ has successive
characters $\psi_0,\ldots,\psi_3$. The factor
$\rho(g_i)-\psi_i(g_i)$ preserves the flag and kills
$F_{i+1}/F_i$. Consequently the ordered product with indices
$0,1,2,3$ is zero: its rightmost factor first lowers $F_4$ to
$F_3$, and the next factors lower successively to zero.
This uses the source order and does not commute the factors.

If $P$ is the polarization matrix, the original form identity
gives
$\rho([g]^*)=P^{-1}\rho(g)^tP$, with
$[g]^*=\epsilon(g)^{1-4}[cgc^{-1}]^{-1}$.
Thus the representation of the star of a killed expression is
also zero, and determinant is unchanged by star. A zero matrix
image of $z$ gives determinant one for $1+h z$ over every
coefficient algebra. Apply the first paragraph to return all
these conclusions to $C$.
\end{proof}
% END JSP2

% BEGIN JSP3
\paragraph{JSP3: the actual ordinary pair and its weight-compatible map.}
\label{claim:JSP3}
Under the retained nonzero-module arithmetic input, there is a
unique $O$-algebra homomorphism
\[
 f:B\longrightarrow\mathbb T,\qquad
 f_*D_B=D_{\mathbb T},\quad f_*\psi_B=\psi_{\mathbb T},\quad
 f\lambda=\text{the structural map }A\longrightarrow\mathbb T.
 \tag{JSP.5}
\]
\begin{proof}
First verify the three relations of the original \texttt{Pair}.
JSP2 applied to the separating polarized arithmetic fibers gives
polarization and the ordered ordinary annihilation identity.
In particular their characteristic coefficients satisfy the
source's star relation, and the trace of $[h]o_v$ is zero.

For the tame relation, put $q=|k|$. Every nonzero residue class
has a unique lift $[a]\in O$ with $[a]^{q-1}=1$.
Indeed successive lifting of a root of $X^{q-1}-1$ works because
its derivative is a unit; completeness gives the limit. If two
lifts coincide residually, factor their difference of
$(q-1)$st powers. The remaining factor is a unit, since its
reduction is $(q-1)a^{q-2}\ne0$, proving uniqueness.
Uniqueness also proves multiplicativity. Thus the residual
characters have the specified Teichm\"uller lifts, and every such
lift takes the value one where its residual character does.
For $g\in I_p(u)=I_u\cap\bigcap_i\ker\overline\chi_i$,
the prescribed tame type makes every specialized $\rho(g)$
conjugate to the identity, hence equal to it. JSP2 gives
$D_{\mathbb T}$ kills $[g]-1$; the required traces are zero.

The remaining \texttt{Pair} fields are exactly the retained
continuity, residual determinant, local-character continuity,
reductions and torsion restrictions. We have therefore constructed
the literal pair over $\mathbb T\in\mathcal C_O$.
Universality of the given $B$ gives $f$ and its uniqueness.
Both $f\lambda$ and the structural map from $A$ carry every
universal inertial character to $\psi_{\mathbb T}|_{I_v}$.
The universal property of the given $\Lambda$ identifies them.
This proves the last identity without an extra choice of weight map.
The Teichm\"uller argument divides by $q-1$, a unit, and not by
$3$ or by $4!$.
\end{proof}
% END JSP3

% BEGIN JSP4
\paragraph{JSP4: recovery of every Hecke generator and surjectivity.}
\label{claim:JSP4}
The map $f$ of (JSP.5) is surjective as an ordinary ring map;
no completion of its image is needed.
\begin{proof}
The algebra $\mathbb T$ is, by its definition, generated over
$A$ by the actual spherical generators $T_{w,j}$
($1\le j\le4$) and ordinary generators $U_{v,j}$
($0\le j<4$). The structural copy of $A$ is in the image
by (JSP.5). Write $q_w=\operatorname{N}w$.
For the original inverse-Frobenius element $g_w$, the source
normalization is
\[
 \chi_{D_{\mathbb T}}([g_w];X)
   =\sum_{j=0}^{4}(-1)^j q_w^{j(j-1)/2}T_{w,j}X^{4-j},
 \qquad T_{w,0}=1.                                  \tag{JSP.6}
\]
Since $w\nmid\ell$, the integer $q_w$ is a unit in $O$.
Taking the coefficient of $X^{4-j}$ gives the explicit preimage
\[
 T_{w,j}=f\!\left((-1)^j q_w^{-j(j-1)/2}
                  [X^{4-j}]\chi_{D_B}([g_w];X)\right).
 \tag{JSP.7}
\]
The sign, inverse Frobenius, and exponent are those of the
original \texttt{heckePol}; none is changed to a different
normalization.

For each $v$, local reciprocity supplies a Weil element $g_v$
whose image under the source Artin map is its specified
uniformizer. The literal uniformizer identity is
\[
 U_{v,j}=\prod_{i=0}^{j}
       \psi_{\mathbb T,v,i}(g_v)\epsilon(g_v)^i
   =f\!\left(\prod_{i=0}^{j}\psi_{B,v,i}(g_v)\epsilon(g_v)^i\right).
 \tag{JSP.8}
\]
Thus every algebra generator is in $f(B)$. Polynomial
expressions in these generators and scalars from $A$ exhaust
$\mathbb T$, proving surjectivity. All inverses in (JSP.7)
are of norms prime to $\ell$, so this proof applies when
$\ell=3$.
\end{proof}
% END JSP4

% BEGIN JSP5
\paragraph{JSP5: the literal ideal and its annihilator are the same.}
\label{claim:JSP5}
If $M=0$, set $J=B$. Otherwise set $J=\ker f$ for (JSP.5).
This is the original \texttt{IsJ} witness. In both cases the
original module has a $B$-action $\mu$ compatible with $\lambda$,
and
\[
               J=\ker\mu=\operatorname{Ann}_B(M).
 \tag{JSP.9}
\]
For $M\ne0$, $B/J\simeq\mathbb T$ as $A$-algebras.
\begin{proof}
In the nonzero case, let $\mu$ be $f$ followed by the inclusion
$\mathbb T\subset\operatorname{End}_A(M)$. This inclusion is
injective because the Hecke algebra was defined as a subalgebra
of endomorphisms. An element of $B$ therefore acts as zero on
$M$ exactly when its image under $f$ is zero. This proves
(JSP.9). The identity $f\lambda=A\to\mathbb T$ gives
$\mu(\lambda(a))m=am$. Surjectivity in JSP4 gives the
quotient isomorphism. The continuous determinant and character
identities in (JSP.5), together with their original
\texttt{IsUniv} formulas, are precisely \texttt{IsMap}.
These are all the nonzero-case \texttt{IsJ} fields.

For $M=0$, its endomorphism ring is the zero ring. The unique
action has kernel $B$ and annihilator $B$; it satisfies the
scalar identity. The source explicitly permits this alternative
$M=0$, $J=B$. We have not inferred a nonzero ordinary module
from a residual datum.
\end{proof}
% END JSP5

% BEGIN JSP6
\paragraph{JSP6: all eight original ring properties.}
\label{claim:JSP6}
The ideal in JSP5 has the complete original
$B.\mathrm{RingProps}(\Lambda,J)$ under the retained arithmetic
input, including the auxiliary unramifiedness assertion stated
there.
\begin{proof}
If $M=0$, then $B/J=0$. It is finite over $A$, its torsion-free
condition is vacuous, and every determinant equality and kernel
identity holds in the zero ring. Suppose henceforth $M\ne0$.
By JSP1, $A$ is a Noetherian domain. Since $M$ is finite free,
$\operatorname{End}_A(M)$ is finite free over $A$.
Its submodule $\mathbb T$ is therefore finite and torsion-free
over $A$. By JSP5 these are exactly \texttt{fin} and
\texttt{tf} for $B/J$ and its specified weight map.

Transport the determinant and characters across
$B/J\simeq\mathbb T$. At a separating arithmetic fiber,
polarization gives the \texttt{polar} field by JSP2.
An upper triangular matrix has characteristic polynomial equal
to the product of its diagonal linear factors. The original
ordered flag therefore gives
\[
 \chi_{D_{B/J}}([\phi_v(g)];X)
       =\prod_{i=0}^{3}(X-\psi_{B/J,v,i}(g)),           \tag{JSP.10}
\]
which is \texttt{O1}. JSP2 gives both the ordinary-product
and star-product kernel assertions \texttt{O2}.

The Teichm\"uller uniqueness proof in JSP3 shows that every
lift allowed by \texttt{IsTeich} is the one already used there.
The tame type of each fiber thus identifies the whole restricted
determinant law on $I_u$ with the diagonal law of these lifts.
JSP2, applied coefficientwise on the group algebra of $I_u$,
proves the original \texttt{types} field, including its
quantifier over all permitted lifts. On $I_p(u)$ those
diagonal matrices are the identity; JSP2 proves both
\texttt{O3} kernel assertions. Finally the retained auxiliary
point statement gives identity matrices on $I_0$ whenever
$r.\mathrm{V0ok}$ holds. The same proof gives both
\texttt{unr0} assertions.

These are the eight source fields: finite generation,
torsion-freeness, polarization, \texttt{O1}, \texttt{O2},
\texttt{O3}, tame type and auxiliary unramifiedness.
The polynomial-law test in JSP2 verifies their assertions over
every coefficient algebra, rather than only at field-valued
points or at traces.
\end{proof}
% END JSP6

% BEGIN JSP7
\paragraph{JSP7: the original weight dictionary and full density.}
\label{claim:JSP7}
For every source strictly increasing weight $h$, the supplied
good Hecke fibers of the corresponding weight imply the exact
$B.\mathrm{Dense}(J)$. In this assertion the weight dictionary
is determined by the original $R.\mathrm{theta}$.
\begin{proof}
Use indices $0,\ldots,3$. At an embedding $\tau$ belonging to
the chosen place $v\mid\ell$, define the dominant weight by
\[
 \lambda_h(\tau)_a=h_{v,3-a,\tau}-(3-a).
 \tag{JSP.11}
\]
Because an increasing integer weight rises by at least one at
each step, these four entries are weakly decreasing.
The pairs $(v,\tau)$ restrict injectively to embeddings of
$L$: the valuation induced on $L$ determines the place, and
density of $L$ in its completion determines the continuous
embedding at that place. Hence (JSP.11) is consistent; assign
zero at unused embeddings. In the source notation
$\mathrm{htW}(\lambda_h,i)=\mathrm{lowW}(\lambda_h,i)+i$,
so it equals $h_{v,i,\tau}$ at the chosen embeddings.
The actual $H$ has $H.\mathrm{admitted}=\mathrm{Set.univ}$,
so every weight $\lambda_h$ just constructed is admitted.

Let $x:B\to\overline{\mathbb Q}_\ell$ and
$z=x\lambda:A\to\overline{\mathbb Q}_\ell$.
For a unit $y$ of $L_v$, let $F_{v,i}(y)$ be the source finite
part
\[
 F_{v,i}(y)=
 \frac{z(\theta(e_{v,i}(y)))}
      {\prod_\tau\tau(y)^{\mathrm{lowW}(\lambda_h,i)(\tau)}}.
\]
For $\sigma$ in the actual inertia group with
$\operatorname{Art}(\sigma)=y$, the \texttt{IsTheta}
identity and $\lambda$'s character identity give
\[
 x(\psi_{B,v,i}(\sigma))F_{v,i}(y)
                  =\prod_\tau\tau(y)^{-h_{v,i,\tau}}. \tag{JSP.12}
\]
Here the explicitly retained source reciprocity normalization has
$\epsilon(\sigma)=\prod_\tau\tau(y)$; its $i$th power
is the shift by $i$ in (JSP.11). This is the local norm formula
for the cyclotomic character in the specified reciprocity;
using the inverse reciprocity map requires inverting both
occurrences, not changing only this shift.

Equation (JSP.12) determines the finite weight-part character
on inertia. The retained uniqueness of the actual Artin map
makes the identity valid for \emph{every} Artin map quantified
in the source's \texttt{IsWtPart}. It is locally constant
exactly when $F_{v,i}$ is
one on a sufficiently deep principal-unit group. Indeed a
locally constant character has open kernel, and reciprocity
identifies its finite quotient on units with a finite quotient
of inertia; conversely such a unit character pulls back to
an open-kernel character on inertia. There are finitely many
$(v,i)$, so the separate depths can be replaced by their maximum.
Thus $x.\mathrm{Arith}(h)$ is exactly the original
$H.\mathrm{Arith}(\theta,z,\lambda_h)$.
Surjectivity of inertia onto units under reciprocity also shows
that the characters for $i\ne j$ are distinct exactly when
$F_{v,i}(y)\ne F_{v,j}(y)$ for some unit $y$.
Consequently $x.\mathrm{Good}(h)$ is exactly
$H.\mathrm{Good}(\theta,z,\lambda_h)$.

Now suppose $b\in B$ vanishes at every good point of weight
$h$ killing $J$. Every good Hecke point
$t:\mathbb T\to\overline{\mathbb Q}_\ell$ yields such
a point $tf$ by (JSP.5) and the dictionary just proved.
It therefore kills $f(b)$. Separation at $\lambda_h$ gives
$f(b)=0$, or $b\in J$. If $M=0$, the conclusion
$b\in J=B$ holds without a point. This is precisely the full
\texttt{Dense} quantifier over every strictly increasing $h$.
\end{proof}
% END JSP7

% BEGIN JSP8
\paragraph{JSP8: the original Taylor--Wiles base on the same ideal.}
\label{claim:JSP8}
The module $M$ and ideal $J$ of JSP5 give the literal
\[
 b:\mathrm{TWBase}(\mathrm{pairOf}(B.u),J,
                         \mathrm{Empty},4[L^+:\mathbb Q]).
 \tag{JSP.13}
\]
Its module is $M$ itself, and its support agrees with $J$
by equality of annihilators, not only by an unproved comparison
of radicals.
\begin{proof}
Let $S=O[[X_1,\ldots,X_d]]$, use JSP1's $e:S\simeq A$,
and set $b.\mathrm{lam}=\lambda e$.
Equip $M$ with its $B$-action $\mu$ and its $S$-action
through $e$. The first identity of (JSP.3) gives
\texttt{wt\_res}; the second, composed with $\lambda$,
gives the exact source variable identity
\[
 b.\mathrm{lam}(X_j)=c_j\psi_{B,v_j,i_j}(\gamma_j)-1.
\]
The scalar identity of JSP5 says the two actions agree along
$b.\mathrm{lam}$. Transport an $A$-basis of $M$ along $e$
to obtain a finite free $S$-basis, including the empty basis
if $M=0$. In that case the zero action and empty basis suffice;
no nonzero Hecke algebra or nonzero-case arithmetic input is needed.
Finally (JSP.9) gives, for every prime $\mathfrak p$,
\[
 \operatorname{Ann}_B(M)\subseteq\mathfrak p
             \quad\Longleftrightarrow\quad J\subseteq\mathfrak p.
\]
These are all fields of \texttt{TWBase}. The replacement
\texttt{pairOf} retains the same underlying determinant and
ordered characters; it changes only the presentation of tame
inertia. The chosen $B$-action and its annihilator are unchanged.
This proves the base construction, not the existence of the
\texttt{TWRing}/\texttt{TWMaps} systems at all later levels.
\end{proof}
% END JSP8

\subsection{The single-point characteristic-prime support argument}

The remaining patching input is not replaced by characteristic-zero
density. We now prove the precise single-point implication used for
\texttt{P1core}, retaining its arithmetic and patching hypotheses.
This separates the elementary changes of rings and the discriminant
argument from the genuinely substantial support theorem.

% BEGIN JSP9
\paragraph{JSP9: topological generation by the original pair values.}
\label{claim:JSP9}
The original universal $B$ satisfies
$\mathrm{GenBy}(O,\mathrm{pairVals}(B.u))$: for every $N$ and
$b\in B$, a polynomial $a$ over $O$ in determinant coefficients
and ordered-character values satisfies
\[
                         b-a\in\mathfrak m_B^N.       \tag{JSP.14}
\]
\begin{proof}
Let $S$ be the set of source pair values. Two $O$-algebra maps
from $B$ to an object of $\mathcal C_O$ agreeing on $S$ give
the same characters and determinant. For the determinant this
uses the integral word identity already used in QX5: its
coefficients on finite sums of group elements are integral
polynomials in characteristic coefficients of group words.
For precision, work modulo any fixed total degree in formal
variables. The unique Lyndon factorization of words gives the
ordered product factorization of
$1-\sum_i Z_i[g_i]$ by factors $1-Z_w[g_w]$ indexed by
primitive cyclic words. Only finitely many factors contribute
at the chosen degree. Applying multiplicativity and homogeneity
of the determinant law expresses its coefficients through
the single-group-element factors. This argument applies to
the law itself and needs no matrix realization of $D_B$.
Thus equality on these coefficients gives equality of the full
laws, also in residue characteristic three. Universality of $B$
then makes the two maps equal.

Put $A_0=O[S]\subset B$ and
$I=(A_0\cap\mathfrak m_B)B$. Suppose
$\mathfrak m_B\ne I+\mathfrak m_B^2$.
The nonzero finite-dimensional $k$-space
$\mathfrak m_B/(I+\mathfrak m_B^2)$ has a nonzero linear
functional. It defines an $O$-derivation $\delta:B\to k$:
subtract from $b$ any lift in $O$ of its residue and apply the
functional. Different lifts differ by an element of
$\mathfrak m_O B\subset I$, so this is well defined.
Addition is immediate, and expanding a product after subtracting
its two scalar residues gives
$\delta(bb')=\bar b\delta(b')+\bar b'\delta(b)$, since
products of elements of $\mathfrak m_B$ are killed.
The two $O$-algebra maps
\[
 b\longmapsto\bar b+\delta(b)\varepsilon,
 \qquad b\longmapsto\bar b
 \quad\text{into }k[\varepsilon]/(\varepsilon^2)
\]
agree on $S$ but are different. This contradicts the preceding
unique-map property. Hence $\mathfrak m_B=I+\mathfrak m_B^2$.

Choose finitely many elements $y_j\in A_0\cap\mathfrak m_B$
whose residues span $\mathfrak m_B/\mathfrak m_B^2$.
Every degree-$r$ graded piece is generated over $k$ by the
degree-$r$ monomials in the $y_j$: express each factor of a
product in $\mathfrak m_B$ modulo $\mathfrak m_B^2$ and
expand. Scalars in $k$ lift to $O$. Starting with a scalar
approximation modulo $\mathfrak m_B$, correct each successive
degree using such a polynomial. After finitely many corrections
this gives (JSP.14). The ring $B$ and the set of character
values have not changed. This is exactly the hypothesis needed
by the source's \texttt{Raw} theorem.
\end{proof}
% END JSP9

% BEGIN JSP10
\paragraph{JSP10: finite coefficient extensions preserve kernels and give stable discriminants.}
\label{claim:JSP10}
Let $j:k_0[[t]]\to k_1[[t]]$ be an injective local homomorphism
between the power-series rings over finite fields occurring in
the source. Its restriction to constants comes from a field
embedding $k_0\to k_1$. In particular
\[
                      \ker(jf)=\ker f.               \tag{JSP.15}
\]
Let $z$ be a normalized source point \texttt{NP} with
$z.a=4$ and let $g:B\to z.K[[t]]$ carry $D_B$ to the
determinant of $z.\rho$. Then every ideal
$\mathfrak q\subseteq\ker g$ has the literal
$\mathrm{DiscSt}(D_B,\mathfrak q)$.
\begin{proof}
For $a\in k_0$, the series $j(a)$ satisfies $X^{|k_0|}=X$.
It is constant: otherwise its least positive nonzero degree
could not occur on the power side, whose corresponding degree
would have to come from a smaller positive degree. Thus $j$
restricts to a field embedding on constants. Locality makes
the constant coefficient of $j(t)$ zero. The induced residue
embedding therefore gives exactly the coefficient action needed
to regard a composite map $jf$ as an $O$-algebra map into the
new power-series ring. Equation (JSP.15) follows from
injectivity and does not require $j(t)=t$.

For the second statement put $F=z.K((t))$.
The normalized span field says that every integral matrix
becomes an $z.K[[t]]$-linear combination of $z.\rho(G')$
after multiplication by some power of $t$; here $G'$ is the
source cyclotomic kernel. Since $z.a=4$, its block algebra
is the full matrix algebra. Inverting $t$ gives
\[
          \operatorname{span}_F z.\rho(G')=M_4(F).
\]
Choose sixteen group images $\rho(g_i)$ forming a basis.
The trace pairing $(U,V)\mapsto\operatorname{tr}(UV)$ is
nondegenerate in every characteristic, since the matrix unit
$E_{ij}$ pairs with $E_{ji}$. Its Gram determinant in this
basis is therefore nonzero. Accordingly
\[
 \Delta=\det\bigl(\operatorname{tr}_{D_B}([g_i g_j])\bigr)_{i,j=1}^{16}
                 \quad\text{has }g(\Delta)\ne0.       \tag{JSP.16}
\]
The determinant specialization identifies each entry with the
matrix trace. If $\mathfrak q\subseteq\ker g$, then
$\Delta\notin\mathfrak q$. Taking the two source lists in
\texttt{DiscSt} to be $(g_i)$ proves the assertion. No division
by the matrix size, a factorial, or the residue characteristic
is involved.
\end{proof}
% END JSP10

\paragraph{The retained support hypotheses.}
For JSP11 retain the source contracts \texttt{Norm},
\texttt{Choice}, \texttt{ArPl},
\texttt{FirstOrderDualArithmeticBound},
\texttt{TaylorWilesObstructionBound}, and \texttt{Raw}, restricted
to the rank-four original data used here. Only their single-point
instances are needed. Their content is as follows.
\begin{enumerate}
\item \texttt{Norm} changes an actual \texttt{NiceI} point, by an
injective finite local extension of its power-series coefficient
ring, to a normalized point. It preserves the determinant and
ordered characters, gives the normalized full-matrix block
$a=4$, preserves \texttt{N3}, and retains \texttt{N4} for the
\emph{original} $D_0$.
\item \texttt{Choice} supplies finitely many elements of $G'$,
one further injective finite local extension, a source cocycle
frame \texttt{Kap}, and distinct eigenvalue lists at those
elements, with the exact \texttt{Detects} property. It
preserves the normalized representation, characters and block
size. The extra comparison of residue-field sizes in the
general source statement is automatic for a one-point family.
\item \texttt{ArPl}, from the given cohomological package
\texttt{CohP}, realizes these elements by a group
$\mathcal T:\texttt{TWGrp}$ with its original
$\texttt{TWPl}$, including all level maps, prime norms,
Frobenius and tame inertia elements, and its subgroup $H$.
It also supplies the source finiteness condition used by
\texttt{Choice}.
\item The two arithmetic bound contracts give exactly
$\mathcal T.\texttt{B1}(k_z,4[L^+:\mathbb Q]+4\mathcal T.q)$
and $\mathcal T.\texttt{B2}(k_z)$ for this normalized point
and its chosen cocycle frame and eigenvalues.
\item \texttt{Raw} gives the source \texttt{ChainConcl} for
the actual pair, ideal and automorphic Taylor--Wiles data,
provided its \texttt{Adapted}, \texttt{GenBy} and point-bound
conditions hold. In the single-point instance $p=\ker g$,
the part used here says: if a prime
$\mathfrak q_1\subseteq p$ has stable discriminant and
contains $J$, then every prime $\mathfrak q\subseteq p$
having stable discriminant contains $J$.
\end{enumerate}
The last assertion is the substantial patched-support input.
The passage to it is not a proof of that input. The exact
quantified source predicates, rather than generic patching
terminology, are bound in the provenance file.

% BEGIN JSP11
\paragraph{JSP11: the complete one-point implication to \texttt{P1core}.}
\label{claim:JSP11}
Fix the same ideal $J\subseteq B$. Suppose that it has the
original $\mathrm{HasAut}(B,J)$, the six retained support
inputs above, and the source side conditions:
$\mathrm{CohP}$, $\mathrm{LocCFT}(s)$, $\mathrm{IsCoef}(\ell,O)$,
$V.S'=s.S$, and a nontrivial residual cyclotomic value on $G_0$.
Then $B.\mathrm{P1core}(J)$ holds. No \texttt{P3core}
assertion or \texttt{NiceE} point is used.
\begin{proof}
Take any original point $x$, map $f:B\to x.K[[t]]$, and
prime $\mathfrak q$ as in \texttt{P1core}: $x$ is
\texttt{NiceI}, $f$ matches the universal pair,
$J\subseteq\ker f$ and $\mathfrak q\subseteq\ker f$.
Use the unchanged arithmetic frame $\mathrm{globOf}(r,V)$.
Its tame inertia is $r.I_p(u)$: an element of the image of
the local inertia lies in the residual kernels precisely when
its local lift lies in their preimages. Consequently
$\mathrm{residOf}(r).I_p(u)=r.I_p(u)$ and this is already
the tame subgroup of the new frame. This verifies the exact
inertia premise for \texttt{Norm}; it also verifies the tame
part of \texttt{Adapted} later.

First take the \texttt{PhiCond} component of \texttt{ArPl}.
Apply \texttt{Norm} to $x$, then the one-point
\texttt{Choice} construction. We obtain a normalized point
$z$ with $z.a=4$, an injective local composite
$j:x.K[[t]]\to z.K[[t]]$, a cocycle frame $k_z$, and the
finite detection and eigenvalue data. JSP10 supplies its
correct $O$-algebra structure and
\[
                         g=jf,\qquad \ker g=\ker f.   \tag{JSP.17}
\]
All determinant and ordered-character identities compose
along these maps, so $g$ still matches the original pair.
The original \texttt{N4} survives the second extension:
it is the vanishing of the kernel of a family of linear
equations in sixteen matrix entries over the fraction field.
A nonzero fraction-field solution could be multiplied by a
common denominator to give a nonzero integral solution, so
the integral \texttt{N4} from \texttt{Norm} gives this assertion.
A finite subsystem has full rank when that kernel is zero,
and its nonzero minor stays nonzero under the injective
field extension. This is the same original auxiliary subgroup,
not the trivial $D_0$ stored in an intermediate frame.

Use \texttt{ArPl} to realize the chosen elements by
$\mathcal T$ and its actual $\mathrm{TWPl}$.
Its identity for $H$ says in particular that every residual
character is one on $H$. Thus the residual characteristic
polynomial of $z.\rho(h)$ is $(X-1)^4$ for $h\in H$;
this follows from the residual determinant identity and the
constant-field embedding of JSP10. We verify the source
$\mathrm{PL}(\ell,H)$ condition used in the two bounds.
If $N=z.\rho(h)-1$, Cayley--Hamilton gives
$N^4\in tM_4(z.K[[t]])$. In characteristic $\ell$,
\[
 z.\rho(h)^{\ell^{4m}}-1=N^{\ell^{4m}}
                                      \in t^mM_4(z.K[[t]]).
 \tag{JSP.18}
\]
The exponent works since $\ell^{4m}\ge4m$.
At each depth $m$ take the open kernel of reduction of the
representation and its normal core in the original profinite
group. The original $\mathcal T.H$ is an open normal subgroup;
its residual-kernel description also makes that normality explicit.
Every conjugate of $h\in H$ is again in $H$, so
(JSP.18) puts a power of $\ell$ of every such $h$ in that
normal core. By the definition of $\mathrm{PL}(\ell,H)$
as the intersection of these eligible open normal subgroups,
the representation is one there modulo $t^m$ for every $m$.
Separation of the power-series ring proves it is exactly one.

The two retained arithmetic bound contracts now apply and
give the precise \texttt{B1} and \texttt{B2} fields printed
above. Along with $z,k_z$, the eigenvalues and $g$, these
construct the source point
$\mathrm{Pt}(\mathcal T,\mathrm{pairOf}(B.u),
4[L^+:\mathbb Q]+4\mathcal T.q)$.
The first \texttt{Adapted} field follows from the same
$H$ identity, and its second field was checked in the first
paragraph. JSP9 supplies \texttt{GenBy} for this exact pair.
The pair replacement preserves all coefficient and character
values, so no new generation argument is required.

The original \texttt{HasAut}(B,J) applies to this very
$\mathcal T$ and its \texttt{TWPl}, supplying
\[
 \mathrm{TWAut}(\mathrm{pairOf}(B.u),J,\mathrm{Empty},
                         4[L^+:\mathbb Q],\mathcal T).
\]
In particular its base ideal
is still $J$, selected before $x$ and $\mathcal T$.
Apply the retained \texttt{Raw} theorem to the one-point
family. Its residue-cardinality comparison is reflexive.
By JSP10, $p=\ker g$ is a prime with stable discriminant,
and every prime below it also has stable discriminant.
Since $J\subseteq p$, choose $\mathfrak q_1=p$ in the
single-point support conclusion. It gives
$J\subseteq\mathfrak q$ for the prime originally chosen,
using (JSP.17). This proves the full original \texttt{P1core}
quantifier. There are no connecting thick primes in a
one-point family and no appeal to their separate
augmentation-kernel comparison.
\end{proof}
% END JSP11

% BEGIN JSP12
\paragraph{JSP12: exact same-ideal assembly and the two-Link quantifiers.}
\label{claim:JSP12}
The preceding proofs reduce the original joint producer to the
following remaining inputs on each of QX4's two actual Links:
the ordinary data in (JSP.2) and the three arithmetic items
following it, together with the printed exact reciprocity matching;
the four remaining Hida fields
\texttt{pt}, \texttt{end1}, \texttt{end2}, \texttt{twist}
on the ideal of JSP5; the all-level automorphic system on this
same ideal; and the retained support inputs of JSP11.
If these inputs are supplied with one common coefficient
threshold after the fixed $\iota$, they give
\texttt{COUNT\_II\_SUPPORT}, hence \texttt{PMExAlt} and
\texttt{CountII} by the existing QX10 and YX2 implications.
\begin{proof}
Fix the original Fit/Deep tower and $\iota$ first.
Choose a coefficient threshold containing QX4's threshold
and any thresholds required by the retained inputs for the
two Links. A finite compositum of coefficient fields gives
one threshold in the original coefficient system. It is
chosen before the later $c'$, before an original Space $X$,
and before its universal data. At every allowed $c'$ use
both QX4 Links with their common $\kappa$ and distinct
auxiliary primes. They keep their own residual data and
universal rings. Work on each with arbitrary actual $B$
and $\Lambda$; no preferred universal representative is
being substituted for the one supplied.

For each Link, JSP3--JSP5 construct its one ideal $J$;
JSP6 and JSP7 give \texttt{ring} and \texttt{dense}.
Combine these with the retained four further fields on
that exact $J$ to obtain \texttt{HidaPropsAlt}.
JSP8 supplies its original Taylor--Wiles base whenever
the retained all-level construction uses that base.
More explicitly, if for every realized $\mathcal T$
the all-level construction supplies the residual lifts
$u_{j,i}$, each $\mathrm{TWRing}_N$, and
$\mathrm{TWMaps}_N$ with one common exponent (the source
construction uses $4!=24$), these records form
$\mathrm{TWAut}$ and hence \texttt{HasAut} on the same $J$.
This observation assembles supplied all-level data; it
does not produce their modules or prove their local operator
identities. If \texttt{HasAut} itself is supplied, JSP11
uses its witness directly.

KX3 supplies the original \texttt{LocCFT} interface.
The original coefficient object gives \texttt{IsCoef}:
it is a complete DVR with finite residue field of
characteristic $\ell$. The Link gives $V.S'=s.S$ and
\texttt{noMu}. The latter means that $L$ does not contain
the $\ell$th roots of unity. Thus the action of $G_L$ on
them is nontrivial. That action factors through the level
quotient, since the omitted inertia is away from $\ell$,
so the required residual cyclotomic value already occurs
on the original $G_0$. Its cohomological input is the
explicit \texttt{CohP} in JSP11; no unnamed strengthening
is added. JSP11 therefore gives \texttt{P1core} on the
very same $J$.

Doing this twice produces $J_0\subset B_0$ and
$J_1\subset B_1$, each with all seven required properties.
Their common threshold precedes all later choices.
QX10 forms
$\phi_0(J_0)X.B+\phi_1(J_1)X.B$ on the original $X$
and gives its five fields. YX2 then supplies the original
\texttt{CountII}. Neither finite composita nor any algebraic
step here imposes $\ell>3$.
\end{proof}
% END JSP12

\paragraph{The remaining boundary and its contribution.}
The kernel and support transport is now a written human argument.
What remains is the production of the ordinary arithmetic system,
the four specified Hida endpoints on that kernel, the actual
all-level automorphic modules, and the six explicitly listed
arithmetic/patching inputs in their single-point scope.
In particular, \texttt{HasAut} is stronger than the
\texttt{TWBase} constructed in JSP8. The latter's finite freeness
and support identity do not construct the higher-level modules or
prove patched support. A rank-two control theorem, ordinary
characteristic-zero density, or QX1's adjoint vanishing does not
fill these missing rank-four inputs.

The frontier remains category C; the exact published/source
matching of a proposed support replacement in characteristic $\ell$
remains category D until its hypotheses and ideal are matched.
The completed JSP implications are category B human arguments,
not new independent existence premises. They are useful on the
RH route because KA5--KA6 remove only Case III for 2-primary
targets: Cases I, II and V still require their actual inputs.
For unrestricted Dirichlet GRH the same \texttt{PMExAlt}
also feeds the selected Case-III route, but it supplies neither
the separate original \texttt{Ctx3} promodularity chain nor
\texttt{SeedIVAlt}. No RH, GRH, or unconditional Count II
conclusion is asserted here.
